\documentclass[11pt]{article}
\pagestyle{empty}

%\setlength{\topmargin}{.25in}

\setlength{\textwidth}{6in}
\setlength{\oddsidemargin}{.25in}
\setlength{\topmargin}{0in}
\setlength{\textheight}{8.25in}


\renewcommand{\baselinestretch}{1.1}
\usepackage{amsmath}
\usepackage{epsfig}
\usepackage{amssymb}
\usepackage{amscd}
\usepackage{times}
\usepackage{graphics}
%\usepackage{chicago}
%\usepackage{fullpage}
%\usepackage{apalike}
\newtheorem{proposition}{Proposition}
\newtheorem{corollary}{Corollary}
\newtheorem{theorem}{Theorem}
\newtheorem{lemma}{Lemma}
\newtheorem{example}{Example}
\newtheorem{claim}{Claim}
\newtheorem{assumption}{Assumption}
\newtheorem{definition}{Definition}
\def\argmax{\mathop{\rm argmax}\limits}
\def\max{\mathop{\rm max}\limits}


\begin{document}

\noindent \large Midterm 2018, ``The Mathematical Foundations of Political Methodology"

\normalsize
\ \\You have 120 minutes for this exam.  You may use a calculator and refer to your Pemberton \& Rau book.  Please do not use the internet during the exam.  Show your work.

\begin{enumerate}

\item Consider the following sets of real numbers:

\begin{eqnarray}
A&=& \{x\ : 1\leq x\leq 5\} \nonumber\\
B&=& \{x\ : 0\leq x\leq 2\}\nonumber\\
C&=& \{x\ : x<0\}\nonumber
\end{eqnarray}

Describe each of the following sets of real numbers in words
\begin{enumerate}
\item $A^c$
%The numbers greater than 5 and less than 1.

\item $A\cup B$
%The numbers between 0 and 5 inclusive.

\item $B\cap C^c$
% B
\item $(A\cup B)\cap C$
% none.
\item $(A\cup B)^c\cap C$
%C

\end{enumerate}

\vspace{.2in}

\item Find the limit of each of the following sequences as $n\rightarrow \infty$, if the limit exists.
\begin{enumerate}
\item $${{4n}\over{(n^2-7)}}$$
% The limit is 0
\item $${{4n^{\frac{1}{2}}}\over{(n^2-7)}}+1$$
% It's 1.
\item  $$\frac{4n}{(n^{\frac{1}{2}}-7)}+1$$
% It doesn't exist.
\end{enumerate}

\vspace{.2in}
\item Use the $(\delta, \epsilon)$ definition of a limit to prove that 
$$\lim_{x\rightarrow 3} 4x-11=1$$
%

\item Use the definitions of continuity and differentiability to evaluate whether or not the following function is 1) continuous and 2) differentiable at $x=0$.
$$f(x)=x^{\frac{1}{3}}$$



\vspace{.2in}
\item Let $$f(x)=\frac{x}{1+e^x}$$
\begin{enumerate}
\item Is f(x) concave, convex or both?
% Take derivative 
\item Find any asymptotes that $f$ may have.
% trace the plot.
\end{enumerate}

\item Perform the following matrix operations:

\begin{enumerate}
\item Find the inverse of $$\begin{bmatrix}1 & 2& 0\\ 0& 3& 1\\1&0&1
\end{bmatrix}$$


\item Find the rank of 
$$A=\begin{bmatrix} 
3& 1& 1&2\\
2& 1& -1&2\\
0& 3& 2&0\\
1& 0& 3&0
\end{bmatrix}$$

%It is 4 

\item What $x$ would make the following matrix singular?
$$\begin{bmatrix}
1&3\\ x&2
\end{bmatrix}$$
%x=2/3
\end{enumerate}

\vspace{.2in}

\item If A and B are matricies, which of the following imply the others: 
\begin{enumerate}
\item AB=BA
\item Either A or B is the identity matrix.
\item A is the inverse of B.
\item A and B are symmetric.
\end{enumerate}
% c implies a, b implies a, a and d imply nothing else.

\item What constant $k$ solves the following equation?
$$\int_0^{10} ky^2dy=1$$

\item What values of the constant $a$ satisfies the following equation?
$$\int_0^a \frac{2x}{1+x^2}dx=1$$
% U = 1+x^2
% du= 2x

%log(1+x^2)|_0^a  
% log(1+a^2)-log(1)=1
% 1+a^2=exp(1)
% a=\sqrt(exp(1)-1)

\end{enumerate}

\end{document}
